%%%
%
% $Autor: Wings $
% $Datum: 2021-05-14 $
% $Pfad: GitLab/MLEdgeComputer/ArduinoScienceJournal.tex
% $Dateiname: 
% $Version: 4620 $
%
% !TeX spellcheck = de_DE/GB
% !TeX program = pdflatex
% !BIB program = biber/bibtex
% !TeX encoding = utf8
%
%%%


\chapter{\TRANS{Arduinos Science Journal}{Arduino's Science Journal}}
    
    
Das Arduino Science Journal ist eine kostenlose und benutzerfreundliche App, die speziell für den Schulunterricht im Bereich Naturwissenschaften entwickelt wurde. Hier sind die wichtigsten Punkte: \cite{}

\begin{enumerate}
  \item \textbf{Zweck:} Das Arduino Science Journal ermöglicht Schülern, wie echte Wissenschaftler zu denken und mit Daten zu arbeiten.
  \item \textbf{Einfach zu verwenden:} Die App ist einfach zu bedienen und kostenlos für Android- und iOS-Geräte verfügbar.
  \item \textbf{Wissenschaftliche Methode:} Schüler lernen die wissenschaftliche Methode kennen und können Experimente durchführen.
  \item \textbf{Sensoren nutzen:} Mit eingebauten Sensoren können Schüler ihre Umgebung erkunden und Daten zu Licht, Bewegung und Klang aufzeichnen.
  \item \textbf{Experimente durchführen:} Die App bietet Dutzende von kostenlosen Experimenten zu verschiedenen Themen wie Physik, Mathematik, Chemie und Biologie.
  \item \textbf{Daten aufzeichnen:} Schüler können Daten aufzeichnen, speichern, exportieren, Diagramme erstellen und Notizen hinzufügen.
  \item \textbf{Teilen und vergleichen:} Experimente können mit anderen geteilt und verglichen werden.
  \item \textbf{Lehrplanabstimmung:} Die App entspricht den Lehrplänen der Next Generation Science Standards (NGSS) in den USA und dem National Curriculum of England.
  \item \textbf{Anpassbarkeit:} Lehrer können die App an verschiedene Lernsituationen anpassen.
  \item \textbf{Lehrerplan:} Der Lehrerplan ist ein Abonnementdienst, der die Integration der App mit Google Classroom ermöglicht.
  \item \textbf{Experimente erstellen:} Lehrer können Experimente direkt in der App erstellen und Aufgaben festlegen.
  \item \textbf{Schülerengagement steigern:} Das neue Design macht die App für Schüler ansprechender.    
  \item \textbf{Effizienz:} Alles an einem Ort - keine Notwendigkeit, zwischen verschiedenen Plattformen zu wechseln.
  \item \textbf{Unterrichtszeit optimieren:} Mehr Zeit für praktisches Lernen im Klassenzimmer.
  \item \textbf{Benutzerbewertungen:} Lehrer und Schüler schätzen die App für ihre praktische Anwendung und Effektivität.
\end{enumerate}

\textcolor{red}{
\begin{enumerate}
    \item Zitieren von NGSS
    \item Beispiele!
    \item \ldots
\end{enumerate}
}


%todo

\URL{https://www.instructables.com/How-to-Use-the-Arduino-Science-Journal-With-the-Na/}
